% ==================================================================
% id:       lim-015
% topic:    גבולות ורציפות של פונקציות
% subtopic: משפטי רציפות (ערך ביניים, ויירשטראס)
% solution: full
% ==================================================================
הוכיחו כי למשוואה
$\ds \lr{x^2-2x+1}^{19}=0$ קיים פתרון בקטע
$(-3,4)$.
\hint{לפני שמנסים את משפט ערך הביניים, שימו לב שהביטוי בסוגריים הוא ריבוע שלם.}
\begin{solbox}
ננסה להשתמש במשפט ערך הביניים. נסמן
$f(x)=(x^2-2x+1)^{19}$.
נשים לב כי
$$,f(-3)=16^{19}>0,\ f(4)=9^{19}>0$$
ובאופן כללי, הפונקציה אי-שלילית, זאת אומרת שכמה שלא ננסה, לא נוכל לספק ערך שלילי לפונקציה, כלומר, לא נוכל להשתמש במשפט ערך הביניים כאן.

ניתן לספק פתרון שלא מסתמך על המשפט הזה.
נשים לב כי
$$f(x)=\lr{x^2-2x+1}^{19}=(x-1)^{38}$$
ו-
$f(1)=(1-1)^{38}=0$.
מכיוון ש-
$1\in(-3,4)$,
השגנו את המטרה.
\end{solbox}
